
Mamba-2 establishes theoretical equivalence between Transformer attention and state space models (SSMs), promising a principled path for knowledge transfer from pretrained Transformers to sub-quadratic architectures \citep{dao2024mamba2}. This structured state-space duality (SSD) framework suggests that attention mechanisms can be converted to SSM form while preserving computational structure---potentially enabling deployment of BERT-scale models with O(n) inference complexity. This work tests whether this duality enables practical distillation and finds a surprising result: duality-derived SSM parameters are numerically stable but reconstruct attention outputs 2.13\% worse than random initialization (p < 0.0001, Cohen's d = $-4.56)$.

This finding is relevant for practitioners considering cross-architecture distillation. If a pretrained Transformer's knowledge could transfer to an efficient SSM via duality equations, edge deployment and long-context inference would become more accessible. However, without understanding the scope and limitations of attention-SSM duality, practitioners may invest compute in approaches that fail at the foundational level.

\subsubsection{The Problem}

Transformers achieve state-of-the-art performance across NLP tasks but suffer from quadratic attention complexity, limiting deployment on resource-constrained devices and long-context applications. Sub-quadratic alternatives---Mamba \citep{gu2023mamba}, RWKV \citep{peng2023rwkv}, RetNet \citep{sun2023retnet}---offer linear-time inference but typically require training from scratch, forfeiting pretrained knowledge.

No validated methodology exists for transferring pretrained Transformer knowledge to SSM architectures while preserving the learned attention structure. Same-architecture distillation (e.g., DistilBERT; Sanh et al., 2019) succeeds but cannot bridge architectural families. Naive output-matching knowledge distillation treats the student as a black box, ignoring the internal computational structure that attention encodes.

The gap addressed here is the empirical validation of Mamba-2 duality for practical distillation. While Mamba-2 proves theoretical equivalence under specific conditions (scalar A matrix, 1-semiseparable causal masks), real attention---multi-head with learned position embeddings---may not satisfy these conditions. Whether duality equations produce SSM parameters that capture meaningful attention structure has not been previously tested.

\subsubsection{Key Insight}

The central finding is that numerical stability does not imply structural fidelity. The duality-to-distillation pipeline is decomposed into two verification stages: (1) parameter validity---do duality equations produce non-divergent SSM parameters?---and (2) structural preservation---do these parameters reconstruct attention outputs better than random initialization?

The first stage passes: 100\% of samples produce stable SSM outputs with 0.$11\times$ magnitude ratio relative to Transformer outputs. The second stage fails: duality initialization is consistently worse than random across all 12 BERT layers. This separation identifies where duality breaks down---at output-level structural fidelity, not at parameter-level validity.

\subsubsection{Contributions}

\begin{enumerate}
\item \textbf{First empirical test of Mamba-2 duality for distillation.} The SSD framework is operationalized for attention-to-SSM parameter conversion and validated on BERT-base, providing evidence about duality's practical scope.
\end{enumerate}

\begin{enumerate}
\item \textbf{Separation of validity from fidelity.} Two-stage verification (h-e1: stability, h-m1: reconstruction) isolates the failure point, showing that valid parameters do not imply preserved structure.
\end{enumerate}

\begin{enumerate}
\item \textbf{Negative result with clear boundary.} Duality initialization produces 2.13\% worse reconstruction error than random (p < 0.0001), with consistent negative effects across all 12 layers. This clarifies when duality should not be expected to provide zero-shot output alignment.
\end{enumerate}

\begin{enumerate}
\item \textbf{Identification of probable causes.} Metric mismatch (output Frobenius norm vs. MOHAWK matrix alignment) and architecture mismatch (multi-head vs. SSD assumptions) are identified as candidates for follow-up investigation.
\end{enumerate}
